% =============================================================================
% 1.7 Bases de données vectorielles et recherche sémantique
% =============================================================================
\section{Bases de données vectorielles et recherche sémantique}
\label{sec:ch1-vecteurs}

\subsection{Principes}

Une base de données vectorielle est un magasin conçu pour stocker des vecteurs
de grande dimension et les interroger par similarité. Elle constitue le
c\oe ur matériel de la recherche sémantique dans un système RAG.

La figure~\ref{fig:base-vectorielle} en rappelle le fonctionnement en deux
temps~: l'indexation, qui projette les documents dans l'espace vectoriel, et
la recherche, qui projette la question dans le même espace puis identifie
les plus proches voisins.

\begin{figure}[htbp]
  \centering
  \resizebox{\linewidth}{!}{% Fonctionnement d'une base de données vectorielle
\begin{tikzpicture}[node distance=5mm and 9mm]

  % --- Indexation ---
  \node[btptitleband, text width=11.5cm] (b1) {Indexation};
  \node[btpbox, text width=2.3cm, below left=6mm and 3.2cm of b1.south] (doc) {Documents};
  \node[btpbox, text width=2.5cm, right=of doc] (emb1) {Modèle de\\plongement};
  \node[btpbox, text width=2.5cm, right=of emb1] (vec1) {Vecteurs\\$\mathbb{R}^{d}$};
  \node[btpdb, right=of vec1, minimum width=2.2cm] (store) {Index\\vectoriel};

  \draw[btparrow] (doc) -- (emb1);
  \draw[btparrow] (emb1) -- (vec1);
  \draw[btparrow] (vec1) -- (store);

  % --- Recherche ---
  \node[btptitleband, text width=11.5cm, below=12mm of doc.south west, anchor=north west] (b2) {Recherche};
  \node[btpboxalt, text width=2.3cm, below=6mm of b2.south west, anchor=north west] (query) {Question};
  \node[btpboxalt, text width=2.5cm, right=of query] (emb2) {Modèle de\\plongement};
  \node[btpboxalt, text width=2.9cm, right=of emb2] (knn) {Plus proches\\voisins\\[-1pt]\btplabel{HNSW, IVF, PQ}};
  \node[btpboxstrong, text width=2.4cm, right=of knn] (rank) {Classement par\\similarité};

  \draw[btparrow] (query) -- (emb2);
  \draw[btparrow] (emb2) -- (knn);
  \draw[btparrow] (knn) -- (rank);
  \draw[btparrowlight] (store.south) -- ++(0,-0.5) -| (knn.north);

  \node[btplabel, below=4mm of knn, text width=6.5cm]
    {Distance cosinus ou euclidienne, calculée de manière approchée
     pour rester sous-linéaire};

\end{tikzpicture}
}
  \caption{Fonctionnement d'une base de données vectorielle}
  \label{fig:base-vectorielle}
\end{figure}

Deux distances dominent la pratique. La \btpterm{similarité cosinus} mesure
l'angle entre deux vecteurs et est invariante à leur norme, ce qui convient
aux plongements dont l'amplitude n'est pas sémantiquement significative. La
\btpterm{distance euclidienne} mesure l'écart géométrique et suppose des
vecteurs comparables en norme. Le choix est dicté par le modèle de
plongement, non par une préférence architecturale.

\subsection{Indexation approximative}

Une recherche exacte des plus proches voisins est linéaire en la taille de
la collection. Pour des corpus de dizaines ou de centaines de milliers
d'extraits, des index approximatifs rendent la recherche sous-linéaire, au
prix d'un rappel légèrement dégradé.

\begin{btptable}{Familles d'index approximatifs}
\label{tab:index-ann}
\begin{tabularx}{\linewidth}{@{}L{2.6cm} Y C{2.2cm} C{2.2cm}@{}}
\toprule
\rowcolor{btppale}
\btpth{Algorithme} & \btpth{Principe} & \btpth{Vitesse} & \btpth{Fidélité} \\
\midrule
HNSW & Graphe hiérarchique de voisinage navigable. & Très élevée & Élevée \\
IVF & Partitionnement en cellules (fichier inversé). & Élevée & Bonne \\
PQ & Quantification produit, compression des vecteurs. & Élevée & Bonne \\
HNSW + IVF & Combinaison graphe et partitionnement. & Très élevée & Très élevée \\
\bottomrule
\end{tabularx}
\end{btptable}

L'algorithme HNSW (Malkov et Yashunin, 2018) est devenu une référence~: il
construit un graphe multi-niveaux dans lequel la recherche descend
progressivement vers le voisinage local de la requête. Sa fidélité élevée et
sa latence prévisible en font un candidat naturel pour une application
interactive.

\subsection{Propriétés recherchées pour les appels d'offres}

Au-delà de la similarité brute, trois propriétés déterminent l'adéquation
d'un magasin vectoriel à notre problématique.

\begin{enumerate}
  \item \btpterm{Filtrage par métadonnées.} La recherche doit pouvoir être
        restreinte à une organisation (cloisonnement), à un type de pièce
        (CCTP, CCAP, référence) et à un rôle (patrimoine ou dossier courant).
        Un index qui ne supporte pas ce filtrage \emph{avant} le calcul de
        similarité oblige à post-filtrer, ce qui dégrade le rappel.
  \item \btpterm{Représentation mixte.} Un magasin capable de stocker
        conjointement un vecteur dense et une représentation creuse permet
        d'exécuter la recherche hybride sans maintenir deux index distincts.
  \item \btpterm{Identifiants stables.} Chaque vecteur doit rester rattachable
        à l'extrait source, à la page, à la pièce et à la version du
        document. Sans cette stabilité, la citation devient impossible.
\end{enumerate}

Le marché propose aujourd'hui plusieurs solutions, ouvertes ou gérées, dont
certaines exposent nativement une recherche hybride ou un modèle de graphe.
Le choix d'un produit particulier relève de l'implémentation
(chapitre~\ref{chap:resultats})~; ce qui importe ici est le cahier des charges
fonctionnel~: filtrage, mixte dense-creux, identifiants citables.
